% GENERATED FILE. Edit canonical lessons, then run npm run generate:books.
% canonical-source-hash: fnv1a64:55237fd799784144
% canonical-lessons: GU-C257-ramat, GU-C257-kheladi, GU-C257-hariphai, GU-C257-bhag-levo, GU-C257-jodavu

\chapter{Games and competitions}
\label{ch:gu-games-and-competitions}
\hlchaptermodality{257}

\begin{chapteropening}
\textbf{By the end of this chapter:} \emph{I can say which game I played and which competition I took part in.}

Complete the last lesson of chapter 257: I can say which game I played and which competition I took part in.
\end{chapteropening}

\section[\texorpdfstring{\gu{રમત} (ramat)}{રમત (ramat)}]{\gu{રમત} (ramat) — game}

\label{lesson:GU-C257-ramat}



\begin{quote}
Before the new one: say the Gujarati for photo, then the Gujarati for to knit.
\end{quote}

\subsection*{You'll want to know: \gu{રમત}}
\textbf{\gu{રમત}} — \emph{ramat} — \textquotedblleft{}game\textquotedblright{}.

\begin{quote}
\textbf{\gu{અમે} \gu{મેદાનમાં} \gu{નવી} \gu{રમત} \gu{રમ્યા}.} — \emph{ame medānmā̃ navī ramat ramyā} — \textquotedblleft{}We played a new game in the field.\textquotedblright{}
\end{quote}

\begin{cousinweb}
\textbf{\gu{રમત}} (\emph{ramat}) is built on the taught \emph{ramvũ}, \textquotedblleft{}to play\textquotedblright{}, with \emph{-at}. CDIAL derives \emph{ramvũ} from Sanskrit \emph{rámatē}, \textquotedblleft{}takes pleasure in\textquotedblright{} (CDIAL 10628).
\end{cousinweb}

\begin{grammarlens}[title={What you've built}]
One more word for everyday life at A2.
\end{grammarlens}

\subsection*{Your turn}
\begin{itemize}
\raggedright
  \item \emph{Say it:} \emph{ramat}
  \item \emph{Say it:} \emph{ramat}, once more
  \item \emph{Say it:} \emph{ame medānmā̃ navī ramat ramyā} — We played a new game in the field.
  \item \emph{Recall:} say the Gujarati for photo, then the Gujarati for to knit, then say \emph{ramat} again
\end{itemize}

\begin{tcolorbox}[breakable,colback=teal!4,colframe=teal!35!black,title={Before you move on}]
What does \gu{રમત} mean? (\textquotedblleft{}Game\textquotedblright{}.) Say it once more.
\end{tcolorbox}
\section[\texorpdfstring{\gu{ખેલાડી} (khelāḍī)}{ખેલાડી (khelāḍī)}]{\gu{ખેલાડી} (khelāḍī) — player}

\label{lesson:GU-C257-kheladi}



\begin{quote}
Before the new one: say the Gujarati for to knit, then the Gujarati for game.
\end{quote}

\subsection*{You'll want to know: \gu{ખેલાડી}}
\textbf{\gu{ખેલાડી}} — \emph{khelāḍī} — \textquotedblleft{}player\textquotedblright{}.

\begin{quote}
\textbf{\gu{તે} \gu{સારો} \gu{ખેલાડી} \gu{હતો}.} — \emph{te sāro khelāḍī hato} — \textquotedblleft{}He was a good player.\textquotedblright{}
\end{quote}

\begin{cousinweb}
\textbf{\gu{ખેલાડી}} (\emph{khelāḍī}) is built on \emph{khel}, \textquotedblleft{}play\textquotedblright{}, with \emph{-āḍī}. CDIAL lists Gujarati \emph{khel} and \emph{khelvũ} under \emph{*khēḍ}, \emph{*khēl-}, \textquotedblleft{}play\textquotedblright{}, reconstructed forms (CDIAL 3918).
\end{cousinweb}

\begin{grammarlens}[title={What you've built}]
One more word for everyday life at A2.
\end{grammarlens}

\subsection*{Your turn}
\begin{itemize}
\raggedright
  \item \emph{Say it:} \emph{khelāḍī}
  \item \emph{Say it:} \emph{khelāḍī}, once more
  \item \emph{Say it:} \emph{te sāro khelāḍī hato} — He was a good player.
  \item \emph{Recall:} say the Gujarati for to knit, then the Gujarati for game, then say \emph{khelāḍī} again
\end{itemize}

\begin{tcolorbox}[breakable,colback=teal!4,colframe=teal!35!black,title={Before you move on}]
What does \gu{ખેલાડી} mean? (\textquotedblleft{}Player\textquotedblright{}.) Say it once more.
\end{tcolorbox}
\section[\texorpdfstring{\gu{હરીફાઈ} (harīphāī)}{હરીફાઈ (harīphāī)}]{\gu{હરીફાઈ} (harīphāī) — competition, contest}

\label{lesson:GU-C257-hariphai}



\begin{quote}
Before the new one: say the Gujarati for game, then the Gujarati for player.
\end{quote}

\subsection*{You'll want to know: \gu{હરીફાઈ}}
\textbf{\gu{હરીફાઈ}} — \emph{harīphāī} — \textquotedblleft{}competition, contest\textquotedblright{}.

\begin{quote}
\textbf{\gu{તે} \gu{હરીફાઈમાં} \gu{જીતી}.} — \emph{te harīphāīmā̃ jītī} — \textquotedblleft{}She won the competition.\textquotedblright{}
\end{quote}

A feminine noun.

\begin{cousinweb}
\textbf{\gu{હરીફાઈ}} (\emph{harīphāī}) is built on \emph{harīph}, \textquotedblleft{}rival\textquotedblright{}, with \emph{-āī}. \emph{Harīph} is a loan from Arabic \emph{ḥarīf}, \textquotedblleft{}a fellow-worker, an associate\textquotedblright{} (Lane), which reached Gujarati through Persian.
\end{cousinweb}

\begin{grammarlens}[title={What you've built}]
One more word for everyday life at A2.
\end{grammarlens}

\subsection*{Your turn}
\begin{itemize}
\raggedright
  \item \emph{Say it:} \emph{harīphāī}
  \item \emph{Say it:} \emph{harīphāī}, once more
  \item \emph{Say it:} \emph{te harīphāīmā̃ jītī} — She won the competition.
  \item \emph{Recall:} say the Gujarati for game, then the Gujarati for player, then say \emph{harīphāī} again
\end{itemize}

\begin{tcolorbox}[breakable,colback=teal!4,colframe=teal!35!black,title={Before you move on}]
What does \gu{હરીફાઈ} mean? (\textquotedblleft{}Competition, contest\textquotedblright{}.) Say it once more.
\end{tcolorbox}
\section[\texorpdfstring{\gu{ભાગ} \gu{લેવો} (bhāg levo)}{ભાગ લેવો (bhāg levo)}]{\gu{ભાગ} \gu{લેવો} (bhāg levo) — to take part}

\label{lesson:GU-C257-bhag-levo}



\begin{quote}
Before the new one: say the Gujarati for player, then the Gujarati for competition.
\end{quote}

\subsection*{You'll want to know: \gu{ભાગ} \gu{લેવો}}
\textbf{\gu{ભાગ} \gu{લેવો}} — \emph{bhāg levo} — \textquotedblleft{}to take part\textquotedblright{}.

\begin{quote}
\textbf{\gu{મેં} \gu{હરીફાઈમાં} \gu{ભાગ} \gu{લીધો}.} — \emph{mẽ harīphāīmā̃ bhāg līdho} — \textquotedblleft{}I took part in the competition.\textquotedblright{}
\end{quote}

Past \textbf{\gu{મેં} \gu{ભાગ} \gu{લીધો}}, \emph{mẽ bhāg līdho}.

\begin{cousinweb}
\textbf{\gu{ભાગ} \gu{લેવો}} (\emph{bhāg levo}), to take a part, is built from \emph{bhāg}, Sanskrit \emph{bhāga}, \textquotedblleft{}portion, fraction\textquotedblright{} (CDIAL 9430; EtymDB), and the verb \emph{levũ}, \textquotedblleft{}to take\textquotedblright{}.
\end{cousinweb}

\begin{grammarlens}[title={What you've built}]
One more word for everyday life at A2.
\end{grammarlens}

\subsection*{Your turn}
\begin{itemize}
\raggedright
  \item \emph{Say it:} \emph{bhāg levo}
  \item \emph{Say it:} \emph{bhāg levo}, once more
  \item \emph{Say it:} \emph{mẽ harīphāīmā̃ bhāg līdho} — I took part in the competition.
  \item \emph{Recall:} say the Gujarati for player, then the Gujarati for competition, then say \emph{bhāg levo} again
\end{itemize}

\begin{tcolorbox}[breakable,colback=teal!4,colframe=teal!35!black,title={Before you move on}]
What does \gu{ભાગ} \gu{લેવો} mean? (\textquotedblleft{}To take part\textquotedblright{}.) Say it once more.
\end{tcolorbox}
\section[\texorpdfstring{\gu{જોડાવું} (joḍāvũ)}{જોડાવું (joḍāvũ)}]{\gu{જોડાવું} (joḍāvũ) — to join}

\label{lesson:GU-C257-jodavu}



\begin{quote}
Before the new one: say the Gujarati for competition, then the Gujarati for to take part.
\end{quote}

\subsection*{You'll want to know: \gu{જોડાવું}}
\textbf{\gu{જોડાવું}} — \emph{joḍāvũ} — \textquotedblleft{}to join\textquotedblright{}.

\begin{quote}
\textbf{\gu{તે} \gu{અમારા} \gu{વર્ગમાં} \gu{જોડાયો}.} — \emph{te amārā vargmā̃ joḍāyo} — \textquotedblleft{}He joined our class.\textquotedblright{}
\end{quote}

\begin{cousinweb}
\textbf{\gu{જોડાવું}} (\emph{joḍāvũ}), to become joined, is built on \emph{joḍvũ}, \textquotedblleft{}to join, fasten\textquotedblright{}, with the intransitive \emph{-ā-}. CDIAL lists \emph{joḍvũ} under \emph{*yōṭayati}, \textquotedblleft{}joins\textquotedblright{}, a reconstructed form in its entry \emph{*yuṭati}, \textquotedblleft{}is joined\textquotedblright{}, and calls the history of the group obscure (CDIAL 10496).
\end{cousinweb}

\begin{grammarlens}[title={What you've built}]
One more word for everyday life at A2.
\end{grammarlens}

\subsection*{Your turn}
\begin{itemize}
\raggedright
  \item \emph{Say it:} \emph{joḍāvũ}
  \item \emph{Say it:} \emph{joḍāvũ}, once more
  \item \emph{Say it:} \emph{te amārā vargmā̃ joḍāyo} — He joined our class.
  \item \emph{Recall:} say the Gujarati for competition, then the Gujarati for to take part, then say \emph{joḍāvũ} again
\end{itemize}

\begin{tcolorbox}[breakable,colback=teal!4,colframe=teal!35!black,title={Before you move on}]
What does \gu{જોડાવું} mean? (\textquotedblleft{}To join\textquotedblright{}.) Say it once more.
\end{tcolorbox}
